{\huge\chapter{Используемые модели и методы}}
\label{sec:Chapter2} \index{Chapter2}

\section{Базовая архитектура}
В работе реализованы и сравниваются две архитектуры: \textbf{CNN-бэкбон (ResNet-50)} и \textbf{Transformer (ViT-Small)}, обе с модулем BNNeck. Изображение \(x\) подаётся в сеть \(f_\theta\), на выходе которой формируется глобальный дескриптор \(z = f_\theta(x)\in\mathbb{R}^d\). Для ResNet-50 используется размер \(256\times128\), для ViT-Small — \(224\times224\). Узел BNNeck нормализует представление перед функциями потерь и в инференсе: \( \hat z = \mathrm{BN}(z) / \lVert \mathrm{BN}(z)\rVert_2 \). % \cite{Luo2019StrongBaseline}

\section{Функции потерь}
\paragraph{Кросс‑энтропия с ослаблением меток.} Классификационная ветвь рассматривает каждый идентификатор как класс; используется Label Smoothing для повышения обобщающей способности.
\paragraph{Triplet Loss с margin.} Для триплетов \((a,p,n)\) (якорь, позитив, негатив) оптимизируется
\begin{equation}
\mathcal{L}_{triplet} = \max\left(0, m + D(f_\theta(a), f_\theta(p)) - D(f_\theta(a), f_\theta(n))\right),
\end{equation}
где \(m>0\) — зазор, \(D(\cdot,\cdot)\) — евклидово или косинусное расстояние. Эффективность повышается за счёт PK‑семплинга (P идентичностей, K изображений каждой) и hard‑mining в батче.

\section{Обучение и аугментации}
Применяются стандартные практики: случайные горизонтальные отражения, Random Erasing, Color Jitter, нормализация по статистикам ImageNet. Оптимизатор Adam/SGD, шаговый или cosine‑annealing scheduler; размер батча подбирается из соображений памяти GPU. Для ускорения сходимости используется заморозка ранних слоёв и прогрев learning rate.

\section{Сопоставление и постобработка}
В инференсе сравнение выполняется по косинусной близости/евклидову расстоянию между L2‑нормированными эмбеддингами. Для улучшения ранжирования применяется k‑reciprocal \emph{re‑ranking}, учитывающий взаимность ближайших соседей и локальную структуру в пространстве признаков. % \cite{Zhong2017Rerank}

\section{Vision Transformer (ViT)}
В качестве альтернативной архитектуры исследуется Vision Transformer (ViT) — модель на основе механизма внимания, изначально разработанная для обработки естественного языка и адаптированная для задач компьютерного зрения.

\paragraph{Архитектура ViT.} Входное изображение разбивается на непересекающиеся патчи размером \(P \times P\) (в работе \(P=16\)). Каждый патч линейно проецируется в эмбеддинг фиксированной размерности \(d\). К последовательности патч-эмбеддингов добавляется специальный токен \texttt{[CLS]}, представление которого после прохождения через Transformer используется как глобальный дескриптор изображения.

\paragraph{Механизм внимания.} В отличие от CNN, где рецептивное поле растёт постепенно, ViT с первого слоя имеет доступ ко всему изображению через self-attention:
\begin{equation}
\text{Attention}(Q, K, V) = \text{softmax}\left(\frac{QK^T}{\sqrt{d_k}}\right)V,
\end{equation}
где \(Q, K, V\) — матрицы запросов, ключей и значений, \(d_k\) — размерность ключей.

\paragraph{Преимущества ViT для Re-ID:}
\begin{itemize}
  \item \textbf{Глобальное внимание:} способность моделировать дальние зависимости между частями изображения с первого слоя;
  \item \textbf{Отсутствие inductive bias:} меньше предположений о структуре данных, что может улучшить обобщение при достаточном количестве данных;
  \item \textbf{Современная архитектура:} активно развивается, множество pretrained моделей.
\end{itemize}

\paragraph{Недостатки ViT:}
\begin{itemize}
  \item \textbf{Требовательность к данным:} без pretrained весов требует больших датасетов;
  \item \textbf{Вычислительная сложность:} квадратичная по числу патчей (\(O(n^2)\));
  \item \textbf{Размер модели:} ViT-Base содержит $\sim$86M параметров vs $\sim$25M у ResNet-50.
\end{itemize}

В работе используется ViT-Small (\(d=384\), $\sim$22M параметров), обученный на ImageNet-21k, с добавлением BNNeck для совместимости с метрическим обучением.

\section{Сравнение архитектур CNN vs Transformer}

\begin{table}[h!]
\centering
\caption{Сравнение архитектур ResNet-50 и ViT-Small}
\label{tab:arch_comparison}
\begin{tabular}{l|c|c}
\hline
\textbf{Характеристика} & \textbf{ResNet-50} & \textbf{ViT-Small} \\
\hline
Тип архитектуры & CNN & Transformer \\
Параметры & $\sim$25M & $\sim$22M \\
Embedding dim & 2048 & 384 \\
Рецептивное поле & Локальное $\rightarrow$ глобальное & Глобальное \\
Inductive bias & Трансляционная инвариантность & Минимальный \\
Pretrain данные & ImageNet-1K & ImageNet-21K \\
\hline
\end{tabular}
\end{table}

Ключевое различие: CNN строят иерархию признаков от локальных к глобальным, тогда как Transformer с первого слоя работает с глобальным контекстом. Для Vehicle Re-ID это потенциально важно, так как распознавание может зависеть как от локальных деталей (фары, решётка радиатора), так и от общей формы кузова.

\section{Детекция транспортных средств: YOLOv11}

В реальных системах видеонаблюдения на вход поступают кадры с нескольких камер, на каждом из которых может присутствовать множество транспортных средств. Перед подачей кадра в модель реидентификации необходима детекция~--- локализация и выделение области интереса (bbox), содержащей отдельное транспортное средство.

В работе используется \textbf{YOLOv11}~--- версия семейства YOLO (You Only Look Once), выпущенная Ultralytics в 2024 году~\cite{ultralytics2024yolo11} и применяемая в данной работе в качестве детектора транспортных средств. Архитектура YOLOv11 относится к классу однопроходных детекторов: за один прямой проход сети одновременно предсказываются координаты рамок и вероятности классов для всех объектов на изображении.

\paragraph{Ключевые характеристики YOLOv11n.} В работе применяется облегчённая версия YOLOv11n ($\approx$2.6M параметров), обеспечивающая баланс между скоростью ($\approx$100~fps на GPU) и точностью (mAP\textsubscript{50} = 39.5\% на COCO). Детектор работает без дополнительного дообучения (zero-shot) на стандартных классах COCO, к которым относятся: \texttt{car} (класс~2), \texttt{bus} (класс~5), \texttt{truck} (класс~7).

\paragraph{Полный пайплайн реидентификации.}
\begin{enumerate}
  \item Входное изображение подаётся в YOLOv11n; детектируются все транспортные средства с порогом уверенности $\geq 0.3$.
  \item Выбирается bbox с наибольшей уверенностью (если детекций несколько).
  \item Кроп изменяется до размера модели ReID ($256\times128$ для ResNet-50, $224\times224$ для ViT-Small).
  \item Кроп подаётся в ReID-модель, извлекается L2-нормированный эмбеддинг.
  \item Эмбеддинги сравниваются косинусным расстоянием с галереей для ранжирования кандидатов.
\end{enumerate}

Если детектор не обнаружил транспортное средство (например, из-за плохого качества изображения), в качестве входа используется полное исходное изображение (fallback).

